\section{Menaces à la validité : synthèse}\label{sec:menaces-ref}

Ce chapitre regroupe les menaces relevées dans les chapitres précédents, classées selon leur effet sur ce que le
projet affirme. Chacune renvoie à sa section détaillée.

\subsection{Sur la précision annoncée}
\begin{itemize}
\item \textbf{Petit échantillon temporel.} Les séries couvrent 2012 à 2026, soit au plus 15 années. Le banc
  départemental compte 22 cellules profession × origine. Les intervalles de confiance sont obtenus par bootstrap des
  départements, ce qui ignore la dépendance entre professions d'un même département et entre origines voisines. Les
  intervalles annoncés sont donc probablement trop étroits\derive.
\item \textbf{Dépendance au millésime.} Une révision de la table privée des sages-femmes a déplacé la MAPE de C1
  de 7,73 à 7,76\,\%\mesure, et celle d'A2 de 8,57 à 8,60\,\%\mesure. Toute comparaison entre deux runs doit être
  faite sur le même millésime (section~\ref{sec:methodes-ref}).
\item \textbf{Bruit incompressible.} Le budget d'erreur attribue 3,11 points\mesure de MAPE au bruit de comptage.
  Un modèle plus riche ne peut pas réduire cette part. Des gains annoncés au-delà de cette borne signaleraient un
  surapprentissage.
\item \textbf{Horizon.} La persistance bat A2 aux horizons 8 et 10 ans\mesure. Les sorties au-delà de dix ans ne
  sont pas des prévisions.
\end{itemize}

\subsection{Sur le pré-enregistrement}
\begin{itemize}
\item Pour 6 protocoles sur 12, l'historique git ne prouve pas que le gel du protocole a précédé le calcul. De plus,
  l'historique local a été réécrit au moins une fois (section~\ref{sec:validation-ref}).
\item Les empreintes recalculées des protocoles gelés concordent toutes\mesure, ce qui garantit leur intégrité
  depuis le gel, mais pas leur antériorité.
\item Plusieurs tests ont été lancés sur les mêmes données. Aucune correction de multiplicité n'est appliquée
  globalement ; seules certaines comparaisons utilisent Bonferroni.
\end{itemize}

\subsection{Sur les mécanismes et les leviers}
\begin{itemize}
\item \textbf{Aucun levier validé de bout en bout.} Le test en différence de différences sur les infirmiers a
  échoué deux fois. Les effets de levier affichés sont des conséquences du modèle, pas des effets mesurés.
\item \textbf{Paramètres supposés.} Une partie des hyperparamètres reste supposée ou introuvable après la relecture
  intégrale des sources (section~\ref{sec:parametres-ref}). Le paramètre $\theta = 0{,}80$ est utilisé par défaut,
  alors que son protocole le déclare non interprétable.
\item \textbf{Pas de rétroaction.} L'effectif ne réagit ni au besoin ni à la densité : l'effet de la densité sur les
  installations mesuré par le protocole N9 est nul. Un transfert de tâches ou une politique d'attractivité ne
  modifie donc pas l'effectif projeté.
\end{itemize}

\subsection{Sur les données}
\begin{itemize}
\item \textbf{Définition du secteur privé.} La séparation public/privé repose sur des fourchettes, pas sur un
  comptage direct ; la largeur de la fourchette dépend de la profession (section~\ref{sec:donnees-ref}).
\item \textbf{Petits effectifs.} Au grain communal, la plupart des cellules comptent moins de cinq personnes : le
  signal y est dominé par le bruit, et ces cellules posent un problème de confidentialité
  (section~\ref{sec:conformite-ref}).
\item \textbf{Ruptures de série.} Les ruptures dans les sources (par exemple la série BTS des infirmiers en 2022)
  et les changements de millésime peuvent créer des sauts que le modèle prend pour des tendances.
\end{itemize}

\subsection{Sur la couche de langage et l'assistant}
\begin{itemize}
\item Le score du compilateur sur le seul jeu resté entièrement à l'aveugle (70\,\%\mesure) est inférieur au score
  agrégé affiché ailleurs (77\,\%) ; les jeux précédents ont servi au réglage.
\item L'apport du LLM face aux réponses déterministes n'a jamais été mesuré. Un cas observé montre l'assistant
  affirmant un effet nul qui n'a pas été estimé (section~\ref{sec:produit-ref}).
\end{itemize}

\subsection{Sur l'évaluation du projet lui-même}
Les notes de la grille du jury (64/100) sont des estimations produites par des agents qui jouaient les jurés. Ce ne
sont pas des mesures ; seuls les contrôles qu'elles citent (latences, recette, compilateur) sont mesurés.
